\documentclass[11pt,a4paper]{article}

\usepackage[margin=2.5cm]{geometry}
\usepackage{amsmath,amssymb,amsthm}
\usepackage{booktabs}
\usepackage{array}
\usepackage{xcolor}
\usepackage{hyperref}
\usepackage{url}
\usepackage{graphicx}
\usepackage{microtype}
\usepackage{parskip}
\usepackage{natbib}
\usepackage{enumitem}
\usepackage{caption}
\usepackage{subcaption}
\usepackage{listings}
\usepackage{algorithm}
\usepackage{algpseudocode}

\hypersetup{
  colorlinks=true,
  linkcolor=blue!60!black,
  citecolor=blue!60!black,
  urlcolor=blue!60!black
}

\lstset{
  basicstyle=\small\ttfamily,
  breaklines=true,
  frame=single,
  language=Python
}

\newtheorem{observation}{Observation}
\newtheorem{definition}{Definition}
\newtheorem{remark}{Remark}

\title{\textbf{An Empirical Operator Algebra on a Latent Manifold\\
       of Polynomial Continued Fractions}\\[0.4em]
       {\large Technical Report / Research Note --- v0.1}}

\author{David Svensson\\
        \small Independent Research\\
        \small \url{https://github.com/VesterlundCoder/cf-space}}

\date{May 2026}

% ──────────────────────────────────────────────────────────────────────────────
\begin{document}
\maketitle

\begin{abstract}
We construct a finite-dimensional coefficient space of polynomial
continued-fraction (PCF) generators and train a Geometric Operator
Transformer Autoencoder (GOT~v3) to embed these objects into a
low-dimensional latent manifold $z \in \mathbb{R}^3$.  We then define
empirical operator vectors $v_T(x) = E(Tx) - E(x)$ for ten symbolic
transformations including index shifts, coefficient scaling, sign flips,
and Ap\'ery-like perturbations.  Using this representation, we test 35
algebraic properties---closure, inverses, involutions, commutativity,
commutator norms, cosine similarity, approximate additivity, Jacobi-type
identities, Lie-bracket closure, braid relations, and conservation
laws---at three levels: symbolic (coefficient space), latent ($z$-space),
and arithmetic (prediction heads).  The results suggest that smooth
operators admit approximate linearisation in latent space, while the
\texttt{shift} operator behaves as a disruptive non-commuting generator.
We identify a small subset of smooth operators with approximate
Lie-algebra-like closure under empirical brackets.  Cosine similarity
$\cos(v_{\mathrm{apery\_pos}}, v_{\mathrm{apery\_neg}}) = -0.998$ confirms
that the Ap\'ery perturbation pair forms the cleanest algebraic inverse in
this space.  These findings are preliminary and checkpoint-dependent, but
demonstrate a reproducible framework for studying operator geometry in
spaces of mathematical objects.
\end{abstract}

\tableofcontents
\bigskip

% ──────────────────────────────────────────────────────────────────────────────
\section{Introduction}

Polynomial continued fractions (PCFs) are mathematical objects of the form
\begin{equation}
  \mathrm{CF}(a, b) \;=\;
  a_0 + \cfrac{b_1}{a_1 + \cfrac{b_2}{a_2 + \cfrac{b_3}{\ddots}}}
\end{equation}
where $a(n) = \sum_{i} a_i n^i$ and $b(n) = \sum_{i} b_i n^i$ are polynomials
with rational coefficients.  PCFs have been central to the study of
irrationality proofs (Ap\'ery's proof of the irrationality of $\zeta(3)$
\cite{apery1979}), fast convergence representations of mathematical constants,
and constructive approaches to hypergeometric identities
\cite{gauss1812hypergeometric}.

A natural question is: \emph{what algebraic structure does the space of PCF
generators possess under symbolic transformations?}  If we think of a
generator $x = (a(n), b(n))$ as a point in a coefficient space
$\mathcal{M}_{\mathrm{CF}}$, then operators such as index shift, scaling, or
sign flip are transformations on this space.  The question becomes: do these
operators form a group, a monoid, a Lie algebra, or something more irregular?

Rather than attempting a purely symbolic algebraic analysis---which would
require enumerating all PCFs---we take a \emph{computational geometry}
approach.  We train a neural autoencoder to embed PCFs into a low-dimensional
Euclidean space $z \in \mathbb{R}^k$, and then measure the algebraic
properties of the resulting operator vectors $v_T(x) = E(Tx) - E(x)$
empirically over a large corpus.

This is related to the broad program of geometric deep learning
\cite{vaswani2017} and the use of latent spaces for structural discovery in
mathematics.  Our specific contribution is a systematic empirical test suite
for 35 algebraic laws applied to a corpus of 286{,}310 PCF generators, with
all code, data, and results publicly available \cite{cfspace2026github}.

\paragraph{Claims.}
We do not claim to have discovered new continued fraction identities.
We do not claim that $\mathcal{M}_{\mathrm{CF}}$ is a Lie algebra.
We claim to present a reproducible empirical framework and to report which
algebraic properties hold approximately and which do not, with explicit error
thresholds.

% ──────────────────────────────────────────────────────────────────────────────
\section{Construction of the CF Space}
\label{sec:dataset}

\subsection{Polynomial continued fraction generators}

Each generator is a pair $x = (a, b)$ where
\begin{align}
  a(n) &= \sum_{i=0}^{d_a} a_i\, n^i, \qquad d_a \in \{1,2,3,4\}\\
  b(n) &= \sum_{i=0}^{d_b} b_i\, n^i, \qquad d_b \in \{1,2,3,4\}
\end{align}
with coefficients $a_i, b_i \in \{-3, -2, \ldots, 3\} \setminus \{0\}$
sampled uniformly.  Each CF is evaluated to depth~400 using \texttt{mpmath}
(50 decimal places).  We retain generators satisfying
\[
  |{\rm CF}_{400} - {\rm CF}_{200}| < 10^{-6}
\]
as a convergence filter.

\subsection{Corpus statistics}

The resulting corpus \texttt{cf\_large} contains $N = 286{,}310$ generators.
Each is stored as a coefficient vector $c \in \mathbb{R}^{20}$ (10
coefficients for $a(n)$, 10 for $b(n)$, zero-padded to degree~9), along
with 20 derived numerical features.

Targets computed per generator:
\begin{itemize}[noitemsep]
  \item \textbf{conv\_exponent}: convergence rate $\log|{\rm CF}_d - L| / d$
  \item \textbf{cf\_quality}: fraction of depths where the partial convergent
        stabilises
  \item \textbf{component}: k-means cluster label ($K=20$, from coefficients)
  \item \textbf{plateau}: binary indicator for exponential convergence plateau
\end{itemize}

\subsection{Relation to known constants}

We use proximity to $\zeta(3) \approx 1.202\,056\,903$ as a scalar field
over the manifold.  This is used as an energy landscape for alignment
experiments in Section~\ref{sec:level3}, \emph{not} as an identity claim.

% ──────────────────────────────────────────────────────────────────────────────
\section{Geometric Operator Autoencoder (GOT~v3)}
\label{sec:model}

\subsection{Architecture}

The model separates geometric and predictive representations:
\begin{align}
  h &= \mathrm{TransformerEncoder}(c) \in \mathbb{R}^{128}
    & &\text{(CLS token, full context)} \\
  z &= \mathrm{Bottleneck}(h) \in \mathbb{R}^k
    & &\text{(geometric manifold, } k = 3\text{)} \\
  \hat{c} &= \mathrm{MLP}_{\rm dec}(z)
    & &\text{(reconstruction from } z \text{ only)} \\
  \hat{\delta} &= \mathrm{MLP}_{\rm delta}([h \| z])
    & &\text{(arithmetic head)}
\end{align}
The TransformerEncoder has $d=128$, $L=4$ layers, $h=8$ attention heads,
totalling 885{,}567 parameters.  The bottleneck dimension $k=3$ is chosen to
match the intrinsic dimensionality $\hat{d} \approx 2.9$ estimated by TwoNN
\cite{levina2005intrinsic} on the corpus.

An additional \emph{op-field} network predicts the latent displacement caused
by an operator:
\[
  \Delta z^{\rm pred} = F_{\rm op}(\tilde{z},\, \delta_c)
\]
where $\tilde{z} = z / \|z\|$ and $\delta_c = c' - c$ is the coefficient
difference.

\subsection{Loss functions}

Training uses a composite loss:
\begin{equation}
  \mathcal{L} = \lambda_r \mathcal{L}_{\rm recon}
    + \lambda_n \mathcal{L}_{\rm nbr}
    + \lambda_o \mathcal{L}_{\rm op}
    + \lambda_c \mathcal{L}_{\rm contrast}
    + \lambda_m \mathcal{L}_{\rm mix}
    + \lambda_z \mathcal{L}_{z\text{-reg}}
    + \lambda_v \mathcal{L}_{\rm var}
\end{equation}

\begin{description}[noitemsep]
  \item[$\mathcal{L}_{\rm recon}$] MSE reconstruction of coefficients from $z$.
  \item[$\mathcal{L}_{\rm nbr}$] Neighbourhood preservation: distances in
        $z$-space (unit sphere) should match coefficient distances.
  \item[$\mathcal{L}_{\rm op}$] Operator consistency:
        $\|\Delta z^{\rm pred} - (\tilde{z}' - \tilde{z})\|^2$,
        computed on normalised $z$ to prevent scale explosion.
  \item[$\mathcal{L}_{\rm contrast}$] SimCLR supervised contrastive loss
        \cite{khosla2020supervised} on normalised $z$ with component labels.
  \item[$\mathcal{L}_{\rm mix}$] Cross-entropy on mixture component logits.
  \item[$\mathcal{L}_{z\text{-reg}}$] L2 penalty $\|z\|^2$ keeping scale bounded.
  \item[$\mathcal{L}_{\rm var}$] VICReg-style anti-collapse
        \cite{bardes2022vicreg}: $\text{mean}(\text{relu}(1 - \text{std}_d(z)))$,
        penalising dimensional collapse without unbounded magnitude growth.
\end{description}

\subsection{Training protocol}

Three-phase training on MPS (Apple Silicon) / CUDA:
\begin{enumerate}[noitemsep]
  \item \textbf{Phase 1 — Geometry} (80 epochs, lr = $3\times10^{-4}$):
        $\lambda_r = \lambda_n = 1$, $\lambda_o = 0.1$,
        $\lambda_c = 0.2$, $\lambda_m = 0.2$,
        $\lambda_z = 0.005$, $\lambda_v = 0.1$.
  \item \textbf{Phase 2 — Arithmetic} (40 epochs, frozen encoder):
        $\lambda_\delta = 1$, $\lambda_{\rm conv} = 0.3$,
        $\lambda_{\rm plateau} = 0.5$.
  \item \textbf{Phase 3 — Joint} (40 epochs, lr = $5\times10^{-5}$):
        all heads active, $\lambda_r = \lambda_n = 0.2$, $\lambda_o = 0.5$.
\end{enumerate}

Phase~1 metrics at epoch~64/80: $\text{mix\_acc} = 0.985$,
$\text{recon} = 0.016$, total loss $= 0.830$.

\paragraph{Note on checkpoint.}
The algebra results in Section~\ref{sec:results} were obtained with an
earlier checkpoint exhibiting partial $z$-collapse.  Symbolic-level
measurements are unaffected.  Latent-level measurements are marked
\textbf{[preliminary]} and will be updated in v0.2 after full pipeline
completion.

% ──────────────────────────────────────────────────────────────────────────────
\section{Operator Definitions}
\label{sec:operators}

\begin{definition}[Operator vector]
  For a CF generator $x = (a(n), b(n))$ and operator $T$, the
  \emph{operator vector} is
  \[
    v_T(x) = E(Tx) - E(x) \in \mathbb{R}^k
  \]
  where $E(x) = z$ is the bottleneck encoding.
\end{definition}

Ten operators are tested, listed in Table~\ref{tab:operators}.

\begin{table}[h]
\centering
\caption{CF operators used in the algebra test suite.}
\label{tab:operators}
\begin{tabular}{lll}
\toprule
Name & Definition & Exact inverse \\
\midrule
\texttt{identity}    & $T(a,b) = (a,b)$ & self \\
\texttt{shift}       & $a(n) \mapsto a(n+1)$ via binomial expansion & \texttt{backshift} \\
\texttt{backshift}   & $a(n) \mapsto a(n-1)$ & \texttt{shift} \\
\texttt{scale\_up}   & $(a,b) \mapsto (1.5a,\, 1.5b)$ & \texttt{scale\_down} \\
\texttt{scale\_down} & $(a,b) \mapsto (a/1.5,\, b/1.5)$ & \texttt{scale\_up} \\
\texttt{apery\_pos}  & $a_2 \mapsto a_2 + 0.10$ & \texttt{apery\_neg} \\
\texttt{apery\_neg}  & $a_2 \mapsto a_2 - 0.10$ & \texttt{apery\_pos} \\
\texttt{sign\_flip\_b} & $b(n) \mapsto b(-n)$ & self (involution) \\
\texttt{sign\_flip\_a} & $a(n) \mapsto a(-n)$ & self (involution) \\
\texttt{random\_sm}  & $(a,b) \mapsto (a + \varepsilon_a,\, b + \varepsilon_b)$,
                       $\varepsilon \sim \mathcal{N}(0, 0.05^2)$ & none \\
\bottomrule
\end{tabular}
\end{table}

The \texttt{shift} operator uses exact binomial expansion:
$a(n+1) = \sum_i a_i (n+1)^i = \sum_i a_i \sum_{j\leq i} \binom{i}{j} n^j$,
preserving polynomial degree.  The \texttt{sign\_flip} operators satisfy
$T^2 = I$ exactly.  The \texttt{apery} operators model small perturbations
inspired by Ap\'ery-type recurrences.

% ──────────────────────────────────────────────────────────────────────────────
\section{Algebraic Tests and Results}
\label{sec:results}

We test 35 laws organised into four levels.  For each law and each operator
(or pair), we report mean error, median, std, max, and a pass/fail indicator
relative to a threshold $\tau$ appropriate to the law.  All tests are run on
$n = 500$ uniformly sampled generators; manifold approximation uses $n = 2000$.

\subsection{Level 1 --- Basic Laws}
\label{sec:level1}

\paragraph{Identity and Inverse.}
Table~\ref{tab:inverses} confirms that all designed inverse pairs are exact
(mean error $< 10^{-5}$) at the symbolic level.

\begin{table}[h]
\centering
\caption{Inverse law errors (symbolic level, $n=500$).}
\label{tab:inverses}
\begin{tabular}{lll}
\toprule
Pair & Mean error & Pass? \\
\midrule
shift $\circ$ backshift & $7.7\times10^{-6}$ & \checkmark \\
backshift $\circ$ shift & $7.7\times10^{-6}$ & \checkmark \\
scale\_up $\circ$ scale\_down & $0.0000$ & \checkmark \\
apery\_pos $\circ$ apery\_neg & $0.0000$ & \checkmark \\
\bottomrule
\end{tabular}
\end{table}

\paragraph{Involutions ($T^2 = I$).}
\texttt{sign\_flip\_a} and \texttt{sign\_flip\_b} satisfy $T^2 = I$ exactly
(mean error $= 0$) --- as expected from the definition $a(n) \mapsto a(-n)$.

\paragraph{Idempotence ($T^2 = T$).}
Only \texttt{identity} is strictly idempotent.  Scale, apery, and sign-flip
operators are \emph{approximately} idempotent (mean error $< 0.026$) due to
small coefficient magnitudes relative to perturbation size.
\texttt{shift} fails strongly (mean error $= 4.84$).

\paragraph{Periodicity.}
\texttt{sign\_flip\_a}, \texttt{sign\_flip\_b} have exact period~2.
Scale and apery operators have approximate period~2 (error $< 0.025$).
\texttt{shift} is aperiodic.

\paragraph{Fixed Points.}
CF~\#214 is the approximate fixed point of both scale operators simultaneously
(min error $5\times10^{-4}$), suggesting a CF whose coefficient magnitudes lie
near the normalisation surface.  The apery operators share an approximate fixed
point at CF~\#361 (min error $2\times10^{-4}$).

\subsection{Level 2 --- Commutativity}
\label{sec:level2}

\paragraph{Commutator norms.}
Table~\ref{tab:commutator} gives the full $10\times10$ commutator norm matrix
$\|[T,S]\| = \|E(T(S(x))) - E(S(T(x)))\|$.

\begin{table}[h]
\centering
\caption{Commutator norm matrix $\|[T,S]\|$ (mean over $n=500$). Rows and
columns ordered as in Table~\ref{tab:operators}. Values $>0.1$ bold.}
\label{tab:commutator}
\footnotesize
\setlength{\tabcolsep}{4pt}
\begin{tabular}{l|cccccccccc}
\toprule
 & I & sh & bsh & s$\uparrow$ & s$\downarrow$ & a$+$ & a$-$ & sf$_b$ & sf$_a$ & rnd \\
\midrule
I        & 0    & 0    & 0    & 0    & 0    & 0    & 0    & 0    & 0    & 0.001 \\
shift    & 0    & 0    & 0    & 0    & 0    & 0.003 & 0.004 & \textbf{4.06} & 0.045 & \textbf{0.19} \\
bshift   & 0    & 0    & 0    & 0    & 0    & 0.001 & 0.001 & \textbf{4.50} & 0.045 & 0.012 \\
s$\uparrow$ & 0 & 0    & 0    & 0    & 0    & 0    & 0    & 0    & 0    & 0.002 \\
s$\downarrow$ & 0 & 0  & 0    & 0    & 0    & 0    & 0    & 0    & 0    & 0.001 \\
a$+$     & 0    & 0.003 & 0.001 & 0  & 0    & 0    & 0    & 0    & 0    & 0.001 \\
a$-$     & 0    & 0.004 & 0.001 & 0  & 0    & 0    & 0    & 0    & 0    & 0.001 \\
sf$_b$   & 0    & \textbf{4.06} & \textbf{4.50} & 0 & 0 & 0 & 0 & 0 & 0 & 0.001 \\
sf$_a$   & 0    & 0.045 & 0.045 & 0  & 0    & 0    & 0    & 0    & 0    & 0.001 \\
rnd      & 0.001 & 0.033 & 0.011 & 0.002 & 0.001 & 0.001 & 0.001 & 0.001 & 0.001 & 0.002 \\
\bottomrule
\end{tabular}
\end{table}

\begin{observation}
The only large non-commutativity in the operator set is between
\texttt{sign\_flip\_b} and \texttt{shift}/\texttt{backshift} (commutator
$\approx 4.1$--$4.5$).  All scale, apery, and identity operators commute with
each other.
\end{observation}

\paragraph{Cosine similarity matrix.}
Table~\ref{tab:cosine} shows $\cos(v_T, v_S)$ between mean operator vectors.

\begin{table}[h]
\centering
\caption{Cosine similarity $\cos(v_T,\, v_S)$ between operator directions.
Bold: $|\cos| > 0.5$.}
\label{tab:cosine}
\footnotesize
\setlength{\tabcolsep}{4pt}
\begin{tabular}{l|cccccccccc}
\toprule
 & I & sh & bsh & s$\uparrow$ & s$\downarrow$ & a$+$ & a$-$ & sf$_b$ & sf$_a$ & rnd \\
\midrule
shift      & 0 & \textbf{+1.00} & +0.30 & \textbf{+0.55} & $-0.38$ & $-0.20$ & +0.21 & $-0.43$ & $-0.13$ & 0.00 \\
backshift  & 0 & +0.30 & \textbf{+1.00} & +0.34 & $-0.02$ & +0.26 & $-0.25$ & +0.14 & +0.11 & 0.00 \\
scale\_up  & 0 & \textbf{+0.55} & +0.34 & \textbf{+1.00} & \textbf{$-0.80$} & +0.09 & $-0.08$ & $-0.33$ & $-0.21$ & 0.00 \\
scale\_dn  & 0 & $-0.38$ & $-0.02$ & \textbf{$-0.80$} & \textbf{+1.00} & $-0.07$ & +0.07 & \textbf{+0.55} & +0.26 & 0.00 \\
apery$+$   & 0 & $-0.20$ & +0.26 & +0.09 & $-0.07$ & \textbf{+1.00} & \textbf{$-0.998$} & +0.22 & $-0.13$ & $-0.08$ \\
apery$-$   & 0 & +0.21 & $-0.25$ & $-0.08$ & +0.07 & \textbf{$-0.998$} & \textbf{+1.00} & $-0.22$ & +0.13 & +0.07 \\
sf$_b$     & 0 & $-0.43$ & +0.14 & $-0.33$ & \textbf{+0.55} & +0.22 & $-0.22$ & \textbf{+1.00} & $-0.06$ & 0.00 \\
sf$_a$     & 0 & $-0.13$ & +0.11 & $-0.21$ & +0.26 & $-0.13$ & +0.13 & $-0.06$ & \textbf{+1.00} & 0.03 \\
random     & 0 & 0.00 & 0.00 & 0.00 & 0.00 & $-0.08$ & +0.07 & 0.00 & 0.03 & \textbf{+1.00} \\
\bottomrule
\end{tabular}
\end{table}

\begin{observation}
$\cos(v_{\mathrm{apery\_pos}},\, v_{\mathrm{apery\_neg}}) = -0.998$: the
Ap\'ery perturbation pair forms a near-perfect antipodal pair in latent space.
This makes it the most reliable algebraic inverse in the corpus for
word2vec-style vector arithmetic \cite{mikolov2013word2vec}.
\end{observation}

\subsection{Level 3 --- Latent Geometry}
\label{sec:level3}

\paragraph{Manifold closure.}
Table~\ref{tab:closure} shows the mean distance from $E(T(x))$ to the nearest
corpus point (manifold approximation).

\begin{table}[h]
\centering
\caption{Manifold closure: mean distance to nearest corpus point after applying $T$.
Baseline (identity) = 0.0012.}
\label{tab:closure}
\begin{tabular}{lrrrl}
\toprule
Operator & Mean dist & Ratio to baseline & Stays on manifold? \\
\midrule
identity    & 0.0012 & $1.0\times$    & \checkmark \\
scale\_up   & 0.0031 & $2.5\times$    & \checkmark \\
scale\_down & 0.0012 & $1.0\times$    & \checkmark \\
apery\_pos  & 0.0012 & $1.0\times$    & \checkmark \\
apery\_neg  & 0.0012 & $1.0\times$    & \checkmark \\
sign\_flip\_b & 0.0029 & $2.3\times$ & \checkmark \\
sign\_flip\_a & 0.0016 & $1.3\times$ & \checkmark \\
random\_sm  & 0.0013 & $1.1\times$    & \checkmark \\
backshift   & 0.1407 & $116\times$    & $\times$ \\
\textbf{shift} & \textbf{1.307} & \boldmath$1075\times$ & $\times$ \\
\bottomrule
\end{tabular}
\end{table}

\texttt{shift} takes generators $1075\times$ further from the manifold than
the baseline, confirming that polynomial index-shifting is not a smooth
transformation in this representation.

\paragraph{Eigenoperator alignment.}
We test whether $v_T(x) \approx \lambda z$ (operator vector aligned with
encoding).  Only \texttt{shift} qualifies ($\cos = -0.714$, $\bar\lambda = 1.0$),
suggesting that index shifting is the principal axis of the latent space.

\paragraph{$\zeta(3)$-conditioned alignment.} \textbf{[Preliminary.]}
We compute a contrast vector $d_{\rm core/far}$ between CFs numerically close
to $\zeta(3)$ and CFs far from it, and measure $\cos(v_T,\, d_{\rm core/far})$.
\texttt{scale\_down} is the only operator with positive alignment ($+0.265$),
suggesting that scaling down coefficients moves generators toward the
$\zeta(3)$-proximal region.  This is presented as an observation about the
scalar field, not as a claim of constructing CF identities.

\paragraph{Lipschitz constants.}
Table~\ref{tab:lipschitz} gives Lipschitz constants
$L_T = \|v_T(x) - v_T(y)\| / \|x - y\|$.

\begin{table}[h]
\centering
\caption{Lipschitz constants for each operator.}
\label{tab:lipschitz}
\begin{tabular}{lrrr}
\toprule
Operator & Mean $L$ & Median $L$ & Max $L$ \\
\midrule
apery\_pos   & 0.0036 & 0.0033 & 0.0057 \\
apery\_neg   & 0.0036 & 0.0033 & 0.0056 \\
scale\_down  & 0.0051 & 0.0035 & 0.0255 \\
scale\_up    & 0.0052 & 0.0038 & 0.0208 \\
sign\_flip\_a & 0.0096 & 0.0075 & 0.0149 \\
sign\_flip\_b & 0.0100 & 0.0083 & 0.0145 \\
backshift    & 0.0148 & 0.0025 & 0.0449 \\
\textbf{shift} & \textbf{0.041} & \textbf{0.0084} & \textbf{11.76} \\
\bottomrule
\end{tabular}
\end{table}

Ap\'ery operators are the smoothest (max $L < 0.006$).  \texttt{shift} has
max $L = 11.76$, indicating singular behaviour at specific coefficient
configurations.

\paragraph{Approximate additivity.}
For smooth operators (scale, apery, sign-flip), the additivity error
$\|v_{T \circ S}(x) - (v_T(x) + v_S(x))\| < 0.01$ passes the threshold.
This supports a key property of the space:

\begin{observation}[Latent linearity of smooth operators]
For $T, S \in \{\mathrm{scale}, \mathrm{apery}, \mathrm{sign\_flip}\}$:
\[
  E(T \circ S(x)) \;\approx\; E(T(x)) + E(S(x)) - E(x)
\]
i.e., operator composition corresponds to vector addition in $z$-space, up
to small error.  This is the analogue of word2vec arithmetic for CF operators.
\end{observation}

\subsection{Level 4 --- Advanced Algebraic Structure}
\label{sec:level4}

\paragraph{Jacobi identity.}
We test $\|[T,[S,R]] + [S,[R,T]] + [R,[T,S]]\| \approx 0$.
Table~\ref{tab:jacobi} shows results for four selected triples.

\begin{table}[h]
\centering
\caption{Jacobi identity test (mean norm of triple bracket sum).}
\label{tab:jacobi}
\begin{tabular}{lrl}
\toprule
Triple $(T, S, R)$ & Mean norm & Pass? \\
\midrule
(apery\_pos, apery\_neg, sign\_flip\_b) & $2.0\times10^{-7}$ & \checkmark (near-exact) \\
(shift, backshift, scale\_up)           & $1.8\times10^{-5}$ & \checkmark \\
(shift, scale\_up, apery\_pos)          & $3.2\times10^{-3}$ & \checkmark (approx) \\
(shift, sign\_flip\_b, apery\_pos)      & $4.06$             & $\times$ \\
\bottomrule
\end{tabular}
\end{table}

The Jacobi identity holds for all smooth-operator triples.  It fails only
when \texttt{sign\_flip\_b} is composed with \texttt{shift}, consistent with
their large commutator (Table~\ref{tab:commutator}).

\paragraph{Lie bracket closure.}
We test whether $[T, S] \approx c \cdot v_R$ for some $R$ in the operator
set ($\cos > 0.9$ threshold).  Results in Table~\ref{tab:lie}.

\begin{table}[h]
\centering
\caption{Lie bracket closure: does $[T,S]$ align with a known operator vector?}
\label{tab:lie}
\begin{tabular}{llrl}
\toprule
$[T, S]$ & Best match $R$ & $\cos$ & Closed? \\
\midrule
{[}scale\_down, apery\_pos{]} & apery\_neg & 0.985 & \checkmark \\
{[}scale\_down, apery\_neg{]} & apery\_pos & 0.984 & \checkmark \\
{[}scale\_up, apery\_neg{]}   & apery\_neg & 0.972 & \checkmark \\
{[}scale\_up, apery\_pos{]}   & apery\_pos & 0.972 & \checkmark \\
{[}shift, sign\_flip\_b{]}    & backshift  & 0.595 & $\times$ \\
{[}backshift, apery\_neg{]}   & backshift  & 0.532 & $\times$ \\
\bottomrule
\end{tabular}
\end{table}

\begin{observation}[Approximate Lie-like subalgebra]
The subset $\{\mathrm{scale\_up}, \mathrm{scale\_down},
\mathrm{apery\_pos}, \mathrm{apery\_neg}\}$ exhibits approximate Lie-algebra-like
closure under empirical latent brackets, with cosine similarity $> 0.97$.
This is a preliminary finding, requiring formal verification in v0.2.
\end{observation}

\paragraph{Braid relations.}
The braid relation $T_1 T_2 T_1 = T_2 T_1 T_2$ is tested for four pairs.
Only $(\mathrm{apery\_pos}, \mathrm{sign\_flip\_a})$ satisfies it (mean error
$= 0.015$); all other tested pairs fail, indicating that the operator set does
not generate a Coxeter group.

\paragraph{Conservation laws.}
Table~\ref{tab:conservation} shows what fraction of CFs preserve structural
properties under each operator.

\begin{table}[h]
\centering
\caption{Conservation fractions: fraction of generators preserving property.}
\label{tab:conservation}
\begin{tabular}{lcccc}
\toprule
Operator & deg$(a)$ & deg$(b)$ & sign$(a_0)$ & norm class \\
\midrule
identity, scale, sign\_flip & 100\% & 100\% & 100\% & 100\% \\
apery\_pos / apery\_neg & 36\% & 100\% & 100\% & 100\% \\
backshift & 27\% & 100\% & 77\% & 100\% \\
shift & 29\% & 99\% & 39\% & 100\% \\
random\_sm & 29\% & 24\% & 70\% & 100\% \\
\bottomrule
\end{tabular}
\end{table}

Notably, \emph{all} operators conserve the norm class (100\%): every
transformed generator remains in a recognisable normalisation class.

% ──────────────────────────────────────────────────────────────────────────────
\section{Interpretation}
\label{sec:interpretation}

\subsection{Algebraic structure summary}

Table~\ref{tab:verdict} summarises which algebraic properties are
approximately satisfied by the full operator set and by the smooth
subset.

\begin{table}[h]
\centering
\caption{Algebraic structure verdict.}
\label{tab:verdict}
\begin{tabular}{lll}
\toprule
Property & Full set & Smooth subset \\
\midrule
Identity exists & \checkmark & \checkmark \\
Inverses (designed pairs) & \checkmark & \checkmark \\
Involutions ($T^2=I$) & \checkmark (sign\_flip) & \checkmark \\
Commutativity & $\times$ (mean comm $= 0.18$) & \checkmark \\
Jacobi identity & partial & \checkmark \\
Lie bracket closure & partial & \checkmark ($\cos > 0.97$) \\
Manifold closure & partial & \checkmark \\
Additivity in $z$ & partial & \checkmark \\
Braid relations & $\times$ (mostly) & partial \\
\bottomrule
\end{tabular}
\end{table}

\paragraph{Best-fit structure.}
The full operator set forms a \emph{non-commutative monoid}: closed under
composition, possessing an identity and soft inverses, but not globally
commutative.  The smooth subset
$\{\mathrm{scale}, \mathrm{apery\_pos}, \mathrm{apery\_neg}\}$ exhibits
approximate Lie-algebra-like structure.  The sign-flip operators generate a
structure resembling $\mathbb{Z}/2\mathbb{Z} \times \mathbb{Z}/2\mathbb{Z}$.

\paragraph{The \texttt{shift} operator as disruptive generator.}
\texttt{shift} exhibits several anomalous properties simultaneously: it is
the only eigenoperator, the only operator that takes generators off the
manifold, has the largest Lipschitz constant (max $L = 11.76$), breaks
commutativity with \texttt{sign\_flip\_b}, and is non-additive in $z$-space.
This is consistent with \texttt{shift} being a non-compact translation
generator in the coefficient representation---analogous to a non-unitary
operator in functional analysis.

% ──────────────────────────────────────────────────────────────────────────────
\section{Limitations}
\label{sec:limitations}

We enumerate limitations explicitly to support reproducibility and future work:

\begin{enumerate}[noitemsep]
  \item \textbf{No proved CF identity.} We do not claim to have found new
        continued-fraction representations of $\zeta(3)$ or any other constant.
        The $\zeta(3)$ alignment measurement is a scalar-field observation.
  \item \textbf{Legacy checkpoint for latent measurements.} The checkpoint
        used for Sections~\ref{sec:level3}--\ref{sec:level4} exhibited partial
        $z$-collapse.  Latent-level results are marked \textbf{[preliminary]}.
  \item \textbf{Sample-based manifold approximation.} The manifold is
        approximated by $n=2000$ corpus points; boundary effects and
        sparse regions are not characterised.
  \item \textbf{Operators are chosen, not exhaustive.} The 10 operators
        represent a structured but limited set.  Other transformations
        (e.g.\ contiguous-fraction equivalences, M\"obius transformations)
        are not tested.
  \item \textbf{Closure is model/dataset-relative.} Lie bracket closure is
        measured relative to the trained model and corpus, not universally.
  \item \textbf{High-precision validation pending.} Convergents are computed
        to depth~400 at dps~50; deeper validation using mpmath at dps~200+
        is left to future work.
  \item \textbf{Error thresholds are empirical.} Pass/fail thresholds
        ($\tau = 0.05$ for most laws) are chosen by inspection rather than
        information-theoretic criteria.
\end{enumerate}

% ──────────────────────────────────────────────────────────────────────────────
\section{Reproducibility}
\label{sec:reproducibility}

All code, configuration files, and derived results are available at:

\begin{center}
\url{https://github.com/VesterlundCoder/cf-space}
\end{center}

The repository \cite{cfspace2026github} includes:
\begin{itemize}[noitemsep]
  \item \texttt{got\_v3/} --- full model package (models, losses, dataset,
        operators, training loop, eval)
  \item \texttt{latent\_operator\_algebra.py} --- algebra test suite
  \item \texttt{got\_v3/configs/} --- all YAML training configurations
  \item \texttt{results/operator\_algebra/algebra\_report.json} --- raw results
  \item \texttt{cfspacealgebra.md} --- full tables
  \item \texttt{requirements.txt} --- Python dependencies
  \item \texttt{data/README\_DATA.md} --- dataset manifest and generation
\end{itemize}

To replicate the algebra results:
\begin{lstlisting}[language=bash, caption={Reproducing the algebra test suite.}]
# Install
pip install -r requirements.txt

# Run full test suite (replace with actual checkpoint path)
python3 latent_operator_algebra.py \
    --ckpt ckpt/got_v3_ae_pretrain_k3.pt \
    --data data/cf_large \
    --n 500 \
    --n_manifold 2000 \
    --no_plots \
    --out results/operator_algebra/
\end{lstlisting}

Random seed: \texttt{--seed 42} (default).  Platform: Apple MPS / CUDA / CPU.

% ──────────────────────────────────────────────────────────────────────────────
\section{Conclusion}
\label{sec:conclusion}

We presented a reproducible empirical framework for testing approximate
operator algebra in a learned latent space of polynomial continued fractions.

The main findings are:
\begin{enumerate}[noitemsep]
  \item The full operator set forms a \emph{non-commutative monoid}.
  \item Smooth operators (scale, apery, sign-flip) satisfy approximate latent
        linearity: $E(T \circ S(x)) \approx E(T(x)) + E(S(x)) - E(x)$.
  \item $\{\mathrm{scale}, \mathrm{apery\_pos}, \mathrm{apery\_neg}\}$ exhibit
        approximate Lie-algebra-like closure under empirical brackets
        ($\cos > 0.97$).
  \item $\cos(v_{\mathrm{apery\_pos}}, v_{\mathrm{apery\_neg}}) = -0.998$:
        the cleanest algebraic inverse in the corpus.
  \item \texttt{shift} is a disruptive non-compact generator: eigenoperator,
        manifold-destructive, non-additive, non-commuting with sign-flip.
  \item \texttt{scale\_down} aligns positively with the $\zeta(3)$ scalar
        field (alignment $= +0.265$).
\end{enumerate}

A version~0.2 of this report is planned after completing the three-phase
training pipeline with the improved GOT~v3 checkpoint (VICReg
anti-collapse + normalised operator consistency loss), which addresses
the $z$-collapse limitation of the current results.

\bibliographystyle{plain}
\bibliography{references}

\end{document}
